\documentclass[a4paper]{article}
\usepackage{graphicx}
\usepackage{amsmath,amssymb}
\usepackage{minted}
\usepackage{multirow}
\usepackage{float}
\usepackage{algorithm}
\usepackage{algpseudocode}
\usepackage{todonotes}
\usepackage{forest}
\usepackage{xstring}
\usepackage{xcolor}
\usepackage[most]{tcolorbox}
\usepackage{xparse}
\usepackage{natbib}
\usetikzlibrary{graphs,graphdrawing}
\usegdlibrary{layered}
% ND Mind Mirror default input files.
\input{/home/donkarlo/Dropbox/repo/nd_language_ai_project/src/nd_language_ai/formal/computer/programming/tex/latex/code_clips/external_graphviz.tex}
\input{/home/donkarlo/Dropbox/repo/nd_language_ai_project/src/nd_language_ai/formal/computer/programming/tex/latex/part/control_sequence/kind/semtag/semtag.tex}
\input{/home/donkarlo/Dropbox/repo/nd_language_ai_project/src/nd_language_ai/formal/computer/programming/tex/latex/code_clips/seehere.tex}

\title{Exponential Moving Average (EMA)}
\date{}
\author{Mohammad Rahmani}
\begin{document}
\maketitle

\section{Definition}
An exponential moving average (EMA), also called exponentially weighted moving average (EWMA), is an online smoother that weights recent values more strongly than older ones.

\section{Methodology}
For a parameter $0<\alpha\leq1$, compute
\[
 m_t=\alpha y_t+(1-\alpha)m_{t-1}.
\]
A typical initialization is $m_1=y_1$. Historical weights decay geometrically, so the previous observations never receive a sharp cutoff as in a fixed-window moving average.

One conventional parameterization uses a nominal span $s$:
\[
 \alpha=\frac{2}{s+1}.
\]
The span is a smoothing convention; it is not a hard window size.

\section{Example: Online Orders}
If yesterday's EMA was 100 orders, today's sales are 130 orders, and $\alpha=0.2$, then
\[
 m_t=0.2(130)+0.8(100)=106.
\]
The smooth estimate reacts to the rise but avoids following its full magnitude.

\section{Algorithm}
\begin{algorithm}[H]
\caption{Causal EMA update}
\begin{algorithmic}[1]
 \State $m_1\gets y_1$
 \For{$t=2,\ldots,n$}
    \State $m_t\gets \alpha y_t+(1-\alpha)m_{t-1}$
 \EndFor
 \State \Return $m_1,\ldots,m_n$
\end{algorithmic}
\end{algorithm}

\section{Advantages}
\begin{itemize}
 \item Constant memory and $O(1)$ work per new observation.
 \item Reacts faster than a very long simple moving average.
 \item Appropriate for streaming dashboards.
\end{itemize}
\section{Limitations}
\begin{itemize}
 \item Sensitive to the value of $\alpha$ and abrupt outliers.
 \item Does not explicitly model calendar seasonality.
 \item Does not quantify trend significance.
\end{itemize}
\section{Practical Guidance}
A larger $\alpha$ gives a more responsive, less smooth curve. Distinguish basic EMA smoothing from Holt and Holt--Winters forecasting methods, which explicitly model level, trend, and possibly seasonality.

\bibliographystyle{plainnat}
\bibliography{/home/donkarlo/Dropbox/repo/nd_shared_research/bibliography/bibliography.bib}
\end{document}
